\section*{ЗАКЛЮЧЕНИЕ}
\addcontentsline{toc}{section}{Заключение}

В результате проведенного исследования были сделаны следующие выводы:

1. Развитие сравнительной политологии представляет собой кумулятивный процесс смены и взаимного дополнения методологических парадигм. Переход от описательного формально-легального подхода первой половины XX века к поведенческой и системной революции 1950--1960-х годов расширил географические и функциональные границы исследований. Последующий плюралистический и неоинституциональный этапы позволили преодолеть крайности квантификации, вернув в фокус внимания уникальность исторического контекста в сочетании со строгостью моделирования институтов.

2. Сравнительный метод в политических исследованиях выполняет важнейшую функцию субститута естественного эксперимента. Он позволяет не только упорядочивать и классифицировать факты политической реальности, но и эксплицировать устойчивые причинно-следственные связи, формируя теоретические обобщения среднего уровня.

3. Эффективность компаративистского исследования напрямую зависит от сбалансированного соотношения теории и метода. Теория выступает смысловым ориентиром, формирующим исследовательские гипотезы и критерии отбора переменных, в то время как сравнительный метод защищает политическую науку от ухода в плоскость чисто нормативных умозрительных суждений, укореняя ее в эмпирической реальности. Современная сравнительная политология сохраняет методологическую неоднородность, что позволяет ей гибко реагировать на актуальные вызовы глобализирующегося мира.